
\subsubsection{Research Questions}

\textbf{RQ1:} Do researchers accept AI-generated documentation suggestions at rates sufficient to justify deployment?

\textbf{RQ2:} Does acceptance generalize across dataset types or require domain-specific tuning?

\textbf{RQ3:} Do users engage thoughtfully with suggestions or passively accept content?

\subsubsection{Data Collection}

\textbf{Deployment Context:} 2-week pilot on HuggingFace with 75 volunteer participants. Each participant documented 1-5 datasets, yielding 1,875 suggestion-response pairs.

\textbf{Documentation Template:} Datasheets for Datasets framework, which specifies 50 documentation dimensions across sections including Dataset Description, Composition, Collection Process, Preprocessing, Uses and Applications, Distribution, and Maintenance.

\textbf{Logging:} All interactions logged to \texttt{suggestions_log.json} (suggestion metadata) and \texttt{user_interactions.json} (user responses).

\subsubsection{Statistical Analysis}

Median and mean acceptance rates are reported with 95\% confidence intervals where applicable. For stratified analysis by dataset type, Kruskal-Wallis H-test evaluates whether acceptance rates differ significantly across groups.

For contextual comparison to code assistance benchmarks, the 92\% median is reported alongside the 65-75\% literature range for GitHub Copilot. Different application domains, user populations, and task characteristics preclude direct statistical comparison, but the contrast provides context on acceptance rates achievable for AI-powered assistance.

\subsubsection{Limitations}

\textbf{Self-Selection Bias:} Pilot participants volunteered and are likely more receptive than the general researcher population. The 22-percentage-point margin above threshold (92\% vs. 70\%) provides buffer for population variance.

\textbf{Scope:} This proof-of-concept validates suggestion acceptance but does not measure downstream effects on documentation completeness, quality, or time-to-completion.

\textbf{Temporal:} Single 2-week deployment cannot evaluate longitudinal stability. Novelty effects versus sustained acceptance versus fatigue remain unknown.
